\documentclass{article}
\title{Practice Problem 2.23}
\author{CSAPP}
\usepackage{listings, booktabs, parskip}
\usepackage{solarized-light}
\usepackage{tabularx}
\begin{document}
\maketitle
\noindent \textbf{ Consider the following C function }
\begin{center}
\begin{lstlisting}[language=C]
        #include <stdio.h>
        int fun1(unsigned word){
            return (int) ((word << 24) >> 24);
        }
        int fun2(unsigned word){
            return ((int) word << 24) >> 24;
        }
\end{lstlisting}
\end{center}
\begin{flushleft}
        A. Assume these are executed as a 32-bit program on a machine that 
        uses two's complement arithmetic. Assume also that right shifts of signed values are
        performed arithmetically, while right shifts of unsigned values are performed logically.
\smallbreak
\begin{tabular}{ccc}
\textbf{w}          & \textbf{fun1(w)} & \textbf{fun2(w)} \\
\midrule
\texttt{0x00000076} & \texttt{0x00000076} & \texttt{0x00000076}\\
\texttt{0x87654321} & \texttt{0x00000021} & \texttt{0x00000021}\\
\texttt{0x000000C9} & \texttt{0x000000C9} & \texttt{0xFFFFFFC9}\\
\texttt{0xEDCBA987} & \texttt{0x00000087} & \texttt{0xFFFFFF87}\\
\bottomrule
\end{tabular}
\\
\bigskip
        B. \textbf{fun1} extracts the least significant byte of \texttt{word} and
        zero-extends it to a 32-bit \texttt{int}. Because both shifts are performed on an
        \texttt{unsigned} value, the right shift fills with zeros regardless of the byte's
        top bit. The result is always in the range $[0,\, 255]$, equivalent to casting the
        low byte as \texttt{(unsigned char)}.

        \medskip
        \textbf{fun2} extracts the least significant byte of \texttt{word} and
        sign-extends it to a 32-bit \texttt{int}. The cast to \texttt{int} before the left
        shift means the subsequent arithmetic right shift propagates the byte's most
        significant bit (bit 7) through the upper 24 bits. The result is in the range
        $[-128,\, 127]$, equivalent to casting the low byte as \texttt{(signed char)}.
        The two functions produce identical results only when bit 7 of the low byte is 0
        (i.e.\ the byte value is less than \texttt{0x80}); otherwise \textbf{fun2} yields
        a negative number while \textbf{fun1} yields a positive one.
\end{flushleft}
\end{document}